\chapter{Requirements Appendix}\label{app:requirements}

This appendix provides an example of how to use the requirements tracking system. 
You can define requirements using the `\. requirement` command or creating an inline requirement with `\ inlineReq{R-ID}` and reference them in your text using `\ reqRef{R-ID}`.
Once requirements are defined, you can reference them from anywhere in your document. This is useful for tracing how the system's design and implementation satisfy the initial requirements.

example:

\begin{enumerate}\setlength{\itemsep}{0pt}\setlength{\parskip}{0pt}\setlength{\parsep}{0pt}
\item \textbf{Diagnosis of fire-resilience states \inlineReq{R1}}\footnote{[\rightarrow R\#] refers to the requirement number \# being created at this point. Full description of the requirment is found in the end of each chapter, i.e \S \ref{sec:requirements_level1}.The requirements framework adopted in this work is detailed in Appendix~\ref{app:se_approach} and \S\ref{sec:requirements_level1}.}
\item \textbf{Anticipation of ignitions \inlineReq{R2}}
\item \textbf{Near--real-time identification of hotspots \inlineReq{R3}}
\item \textbf{Near--real-time monitoring of fire-front evolution \inlineReq{R4}}
\item \textbf{All weather, cloud/smoke-penetrating, day/night capability \inlineReq{R5}}
\end{enumerate}


% The descripiton of the `.\\requirement` command has the following structure:
% \requirement{<ID>}{<Title>}{<Description>}{<Type>}{<Tags>}
% - <ID>: A unique identifier for the requirement (e.g., R-01).
% - <Title>: A short, descriptive title.
% - <Description>: A detailed explanation of the requirement.
% - <Type>: The category of the requirement (e.g., Functional, Non-Functional, Performance, Security).
% - <Tags>: Comma-separated tags for filtering or grouping (e.g., UI, Database, API).

\section{Example: Mission Objectives (Chapter 1)}\label{sec:app_req_ch1}
\subsection{Mission Earth Observation Products}

The following requirements define the primary Earth Observation inputs and products necessary to support wildfire resilience and monitoring objectives.

\requirement{R1}{Fire Resilience Inputs}{Provide EO measurement products sufficient for external wildfire-resilience indices; e.g., vegetation state \& dynamics, moisture, thermal regime, and structure; cadence and spectral bands: TBD.}{Mission}{EO product}

\requirement{R2}{Anticipation of ignitions}{Provide sufficient EO capabilities for third-party short-term ignition-risk assessment and forecasts.}{Mission}{EO product}

\requirement{R3}{Near-real-time identification of hotspots}{Provide capability for third-party near-real-time hotspot detection alerts with geo-referenced location and confidence; minimum detectable size, geolocation accuracy, and alert latency: TBD.}{Mission}{EO product}

\requirement{R4}{Near-real-time monitoring of fire-front evolution}{Provide capability for near-real-time third-party products describing the fire front perimeter and spread rate during overpass; update cadence, geolocation accuracy, and revisit: TBD.}{Mission}{EO product}





